\documentclass[a4paper,12pt]{article}
\usepackage[utf8]{inputenc}
\usepackage{amsmath,amsfonts,amssymb}
\usepackage{geometry}
\usepackage{enumitem}

% Pengaturan Margin
\geometry{top=1in, bottom=1in, left=1in, right=1in}

\begin{document}

% --- KOP HEADER MIRIP GAMBAR ---
\begin{center}
    \textbf{SOAL LATIHAN/KUIS/AKHIR} \\
    \textbf{BIMBEL INTENSIF PERSIAPAN OSN-K SMA/MA 2026} \\
    \textbf{18-23 MEI 2026} \\
    \textbf{BIDANG MATEMATIKA} % Garis bawah untuk mengisi bidang
\end{center}

\vspace{0.5cm}

\noindent
\begin{tabular}{lll}
    MINGGU KE-/TANGGAL & : &  1/18 MEI 2026\\
    PUKUL              & : &  11.15-12.00\\
    TOPIK              & : & TEORI BILANGAN \\
\end{tabular}

\vspace{0.4cm}
\hrule height 1.5pt
\vspace{0.8cm}
% --- AKHIR KOP ---

\section*{SOAL UJIAN}

\begin{enumerate}[label=\textbf{\arabic*.}]
    
    % Nomor 1: Simon's Favorite Factoring Trick (SFFT)
    \item Tentukan semua pasangan bilangan bulat positif $(x, y)$ yang memenuhi persamaan $xy = 2x + 2y$. Kemudian, carilah nilai minimum dari $x + y$.

    % Nomor 2: Keterbagian & Digit
    \item Tentukan banyaknya bilangan 3-digit $N$ yang habis dibagi oleh 7, sedemikian rupa sehingga jika urutan digitnya dibalik, bilangan yang baru terbentuk habis dibagi oleh 5.

    % Nomor 3: Teori Bilangan (KPK & FPB)
    \item Diketahui $m$ dan $n$ adalah bilangan bulat positif yang memenuhi persamaan $\text{KPK}(m,n) - \text{FPB}(m,n) = 103$. Tentukan \textbf{jumlah} semua kemungkinan nilai dari $m+n$.

    % Nomor 4: Modulo Faktorial (Digit Terakhir)
    \item Tentukan banyaknya bilangan bulat positif $n \le 2013$ sedemikian rupa sehingga hasil penjumlahan deret faktorial $1! + 2! + 3! + \dots + n!$ memiliki angka satuan (digit terakhir) bernilai 3.

    % Nomor 5: Modulo & Counting
    \item Tentukan banyaknya tripel bilangan bulat positif $(a,b,c)$ dengan batasan nilai masing-masing $1 \le a, b, c \le 5$ sedemikian rupa sehingga hasil kali ketiganya yaitu $abc$ \textbf{tidak habis dibagi} oleh 4.

    % Nomor 6: Faktorisasi Eksponen & Prima
    \item Diketahui bahwa bentuk bilangan $4^{11}+1$ habis dibagi oleh suatu bilangan prima $p$ yang nilainya lebih besar dari 1000. Tentukan nilai dari bilangan prima $p$ tersebut.

    % Nomor 7: Modulo & Sifat FPB
    \item Tentukan banyaknya bilangan bulat positif $n \le 100$ yang memenuhi persamaan $\text{FPB}(n, 2^n-1) = 3$.

\end{enumerate}

\newpage

% --- HALAMAN KUNCI JAWABAN ---
\begin{center}
    \textbf{KUNCI JAWABAN}
\end{center}
\vspace{0.5cm}

\begin{enumerate}[label=\textbf{\arabic*.}]
    \item $8$ 
    \item $14$
    \item $435$
    \item $2011$
    \item $54$
    \item $2113$
    \item $8$
\end{enumerate}

\end{document}